% texte indiquant le nom du cours, formulation comprise
\newcommand{\texteNomCours}{Cours de}  

% texte indiquant le nom du professeur encadrant, formulation comprise
\newcommand{\texteNomProf}{cours dispensé par } 

% numéro du chapitre
\newcommand{\numeroChapitre}{0}

% nom du chapitre
\newcommand{\texteNomChapitre}{} 

\documentclass[12pt]{article}

\usepackage[a4paper,top=2.5cm,bottom=2.5cm,left=1.5cm,right=1.5cm]{geometry}

% Packages de base
\usepackage[utf8]{inputenc}
\usepackage[T1]{fontenc}
\usepackage{fancyhdr}      % pour personnaliser en-têtes et pieds
\usepackage[french]{babel}   % <-- gère la langue

% Maths
\usepackage{amsmath}
\usepackage{amssymb} % pour mathbb
\usepackage{mathtools} 
\usepackage{esint} % intégrale avec un rond
\usepackage{dsfont} % pour l'indicatrice


% Trucs en plus
\usepackage{xcolor} % gestion des couleurs
\usepackage{listings} % pour le code source
\usepackage{enumitem} % pour personnaliser les listes
\usepackage{multicol} % pour les listes en plusieurs colonnes
\usepackage{stmaryrd} % pour les intervalles d'entiers
\usepackage{graphicx} % pour insérer des images
\graphicspath{ {figures/} }

\usepackage[framemethod=tikz]{mdframed} % les boites jolies

% Les graphiques
\usepackage{tikz}
\usepackage{pgfplots}
\pgfplotsset{compat=1.18}


% Style des titres de sections et sous-sections
\usepackage{titlesec}

\usepackage{tocloft} % pour personnaliser la table des matières


\usepackage[colorlinks=true, urlcolor=blue]{hyperref} % les liens

\renewcommand{\thesection}{\Roman{section}} 
\titleformat{\section}
  {\Large\bfseries} % Style du titre
  {\thesection} % Numéro de section
  {0.5cm} % Espacement entre numéro et titre
  {} % Pas de préfixe supplémentaire
\titleformat{\subsection}
  {\large\bfseries} % Style du titre
  {\thesubsection} % Numéro de section
  {0.4cm} % Espacement entre numéro et titre
  {} % Pas de préfixe supplémentaire

\setlength{\cftsecnumwidth}{3em}   % espace réservé au numéro de section
\setlength{\cftsubsecnumwidth}{2.5em} % pour les sous-sections

% Réglages fancyhdr
\pagestyle{fancy}
\fancyhf{} % vide tout


\renewcommand{\headrulewidth}{0.4pt}
\renewcommand{\footrulewidth}{0.4pt}

\fancyfoot[L]{Notes de Raphaël JONTEF}
\fancyfoot[C]{\thepage}
\fancyfoot[R]{Enseirb-Matmeca}
\fancyhead[L]{\texteNomCours}
\fancyhead[R]{C.\numeroChapitre~:~\texteNomChapitre} 

\input{backend/boites_cours.tex}
\input{backend/listings.tex}
\begin{document}
\input{backend/commands.tex}

% Page de titre custom

\hypersetup{pageanchor=false}
\setlength{\headheight}{15pt}
\begin{titlepage}
    \centering
    \vspace*{3cm}
    {\huge Chapitre \numeroChapitre\par}
    {\Huge \bfseries\texteNomChapitre \par} % nom  du chapitre
    \vspace{2cm}
    {\Large Notes rédigées par Raphaël JONTEF\par}
    \vspace{0.5cm}
    {\normalsize \texteNomProf \par}
    \vspace{4cm}
    {\small Enseirb-Matmeca\par}
    {\small filière informatique\par}
    \vfill
    {\large \today\par}
\end{titlepage}
\hypersetup{pageanchor=true} % réactive pour le reste

\tableofcontents
\newpage

\input{sections/introduction}
\section{}
\begin{multicols}{2}


\begin{exercice}{}{}
    Déterminer le signe de chacune des expressions suivantes :
    \begin{enumeratebf}
        \item $A=(-15)+85$
        \item $B=(-15)+(-85)$
        \item $C=15+(-85)$
        \item $D=(+15)+85$
        \item $E=15+(-15)+15+(-15)$
    \end{enumeratebf}
    Calculer alors ces expressions.
\end{exercice}

\begin{exercice}{}{}
    Calculer les sommes suivantes~:
    \begin{enumeratebf}
        \item $A=(-6)+10$
        \item $B=12+(-5)$
        \item $C=(-15)+17$
        \item $D=(-14)+(-23)$
        \item $E=(-25)+(-25)+(-25)+(-25)$
    \end{enumeratebf}
\end{exercice}
\begin{exercice}{}{}
    Calculer les sommes suivantes~:
    \begin{enumeratebf}
        \item $A=(+14)+(-5)$
        \item $B=15+(+158)$
        \item $C=(-47)+78$
        \item $D=(-123)+(-89)$
    \end{enumeratebf}
    \textit{On pourra poser les additions pour les calculer.}
\end{exercice}

\begin{exercice}{}{}
    Recopier et compléter dans les égalités suivantes les "$\dots$" par "$+$" ou "$-$".
    \begin{enumeratebf}
        \item $A = (-20)+17 = \dots(-20 \dots 17)$
        \item $B = 23 + (-10) = \dots(23 \dots 10)$
        \item $C = -53 - 47 = \dots(53 \dots 47)$
        \item $D = (+65) + (+87) = \dots(65 \dots 87)$
    \end{enumeratebf}
\end{exercice}

\end{multicols}
\section{}
\begin{multicols}{2}
    \begin{exercice}{}{}
        Déterminer l'opposé des nombres suivants. Les calculer si nécessaire.
        \begin{enumeratebf}
            \item $-3$
            \item $-5,2$
            \item $-(-7)$
            \item $-(-2,5) + 3$
        \end{enumeratebf}
    \end{exercice}


    \begin{exercice}{}{}
        Effectuer les calculs suivants :
        \begin{enumeratebf}
            \item $A = 4 - (-5)$
            \item $B = -5 - (-10)$
            \item $C = -23 - (+14)$
            \item $D = 100 - (+25)$
        \end{enumeratebf}
    \end{exercice}
    \begin{exercice}{}{}
        Effectuer les calculs suivants :
        \begin{enumeratebf}
            \item $E = -3 - (-7) + 5$
            \item $F = 10 - (-2) - 4$
            \item $G = -15 - (+6) + 8$
            \item $H = 20 - (+10) + 3$
        \end{enumeratebf}
    \end{exercice}
    \begin{exercice}{}{}
        Effectuer les calculs suivants~:
        \begin{enumeratebf}
            \item $A = (-5) - (-3)$
            \item $B = -4,3 - 5,7$
            \item $C = 12,3 - (+4,7)$
            \item $D = 10 - 20,5$
        \end{enumeratebf}
    \end{exercice}
\end{multicols}
\section{}
\begin{multicols}{2}
    \begin{exercice}{}{}
        Sans calcul, déterminer le signe des produits suivants.
        \begin{enumeratebf}
            \item $A=(-154)\times 97$
            \item $B=(-4)\times (-27)\times 103$
            \item $C=42\times (-2)\times (-87)\times (-9)$
            \item $D=(-74)\times 562\times (-2)\times 5\times (-87)$
            \item $E= (-1)^{2026} = (-1)\times \dots \times (-1)$ (2026 facteurs)
        \end{enumeratebf}
        \textit{Il suffit de compter le nombre de facteurs négatifs dans le produit. Si ce nombre est pair, le produit est positif, sinon il est négatif. Par exemple, $F = (-3) \times (-1) \times (-2)$ a 3 facteurs négatifs, donc $F$ est négatif.}
    \end{exercice}

    \begin{exercice}{}{}
        Effectuer les calculs suivants :
        \begin{enumeratebf}
            \item $A = 2,5 \times 2 \times (-1)$
            \item $B = 10 \times (-2) \times (-0,1)$
            \item $C = (-1) \times (-1) \times (-1) \times (-1)$
            \item $D = 4 \times (-2) \times 3 \times 10$
        \end{enumeratebf}
        \textit{On pourra essayer de regrouper certains facteurs pour les calculer à part, puis remplacer l'expression par le résultat. Par exemple, $12 \times 25 \times 4 = 12 \times (25 \times 4) = 12\times 100 = 1200  $}
    \end{exercice}

    \begin{exercice}{}{}
        Effectuer les calculs suivants :
        \begin{enumeratebf}
            \item $A = 3 \times (-4)$
            \item $B = -5,2 \times 6$
            \item $C = -7 \times (-2)$
            \item $D = -2,5 \times (-3) + 4$
        \end{enumeratebf}
    \end{exercice}

    \begin{exercice}{}{}
        Effectuer les calculs suivants, en posant les multiplications :
        \begin{enumeratebf}
            \item $A = 23 \times 5$
            \item $B = 12,5 \times 4$
            \item $C = 7,2 \times 3$
            \item $D = 15 \times 6$ 
        \end{enumeratebf}
    \end{exercice}
\end{multicols}

\section*{}
\input{sections/conclusion}


\end{document}